\section{Capítulo~\ref{sec:sleep}. Sono}

Os modos do sono, os replays e a redução das sinapses foram medidos por registros de neurônios, microdiálise e microscopia eletrônica; o horário do sono e a gangorra lenta são descritos pelos modelos do livro, e a finalidade do sono é, em grande parte, hipótese.

{\footnotesize
\setlength{\tabcolsep}{3pt}\renewcommand{\arraystretch}{1.15}
\begin{longtable}{>{\raggedright}p{0.5cm}>{\raggedright}p{4.0cm}>{\raggedright}p{5.6cm}>{\raggedright}p{2.3cm}>{\raggedright\arraybackslash}p{1.7cm}}
\toprule
N.º & Proposição & Experimento e o que foi medido & Fonte & Status \\
\midrule\endfirsthead
\toprule
N.º & Proposição & Experimento e o que foi medido & Fonte & Status \\
\midrule\endhead
1 & No sono os neurônios do córtex não se calam; o que muda é o modo de funcionamento & Registros longos de neurônios do córtex em ratos: no sono de ondas lentas a taxa continua não nula, e períodos de atividade se alternam com períodos de silêncio; após vigília longa a taxa é maior em todos os estados, após o sono, menor & \cite{Vyazovskiy2009} & medido \\
2 & O \nt{ACET} é alto de dia e no sono REM, baixo no sono de ondas lentas (tabela~\ref{tab:sleepmodes}) & Microdiálise no córtex de gatos em livre movimento: na vigília, a liberação de acetilcolina é $100$--$175\%$ maior que no sono de ondas lentas, e no sono REM fica mais ou menos no nível da vigília tranquila & \cite{Marrosu1995} & medido \\
3 & \nt{NORA} e \nt{SERO} se aquietam no sono de ondas lentas e quase se calam no REM & Registro de neurônios \nt{NORA} isolados em ratos e de neurônios \nt{SERO} em gatos em cada estágio do sono: a taxa é máxima na vigília, menor no sono de ondas lentas e quase nula no REM & \cite{AstonJonesBloom1981,McGintyHarper1976} & medido \\
4 & No sono REM os músculos são desligados por inibição via \nt{GABA} e \nt{GLYC} & Em ratos, mediu-se a atividade dos neurônios \nt{ACET} do músculo masseter e aplicaram-se bloqueadores diretamente sobre eles: a paralisia só cessa se forem bloqueados tanto os receptores \nt{GABA} dos dois tipos quanto os receptores \nt{GLYC} & \cite{BrooksPeever2012} & medido \\
5 & A pressão do sono $S$ é um capacitor com dois limiares~\eqref{eq:sleeppressure}, $\tau_r=18.2$~h, $\tau_d=4.2$~h & Modelo de dois processos; as constantes de tempo foram obtidas de como a potência das ondas lentas do EEG cresce com a vigília e cai no sono, e a forma dos limiares, do horário do despertar espontâneo & \cite{Borbely1982,Daan1984} & teoria \\
6 & A máquina chega sozinha a $16.1$~h de vigília e $8.0$~h de sono (fig.~\ref{fig:twoprocess}) & Cálculo por~\eqref{eq:sleeppressure} com limiares oscilantes, script \texttt{models/sleep\_models.py}: $16.1$~h de vigília e $7.9$--$8.0$~h de sono & --- & modelo \\
7 & Após $40$~h sem dormir, $S=0.90$, mas o sono dura só $8.8$~h: a dívida é paga com profundidade, não com duração & Modelo: \texttt{models/sleep\_models.py} dá $S=0.90$ e $8.8$~h. Experimento: após $40.5$~h sem dormir, em oito pessoas, na primeira noite de recuperação, a potência das ondas lentas do EEG foi significativamente maior que o normal & \cite{Borbely1981} & modelo \\
8 & Fisicamente, $S$ é a adenosina & Microdiálise em gatos: a adenosina extracelular numa das regiões que controlam a vigília cresce durante a vigília, cresce ainda mais com a privação de sono e cai durante o sono. Que $S$ seja só adenosina não foi mostrado & \cite{PorkkaHeiskanen1997,Dunwiddie2001} & hipótese \\
9 & No sono de ondas lentas o córtex oscila: atividade por algumas centenas de ms, depois silêncio, menos de uma vez por segundo ($<1$~Hz, sob anestesia em geral $0.3$--$0.4$~Hz) & Registro intracelular em gatos anestesiados: despolarização de $0.8$--$1.5$~s e hiperpolarização longa, ritmo $<1$~Hz, em geral $0.3$--$0.4$~Hz; a mesma alternância de atividade e silêncio aparece no sono natural (linha~1) & \cite{Steriade1993,Vyazovskiy2009} & medido \\
10 & A gangorra é um grupo com realimentação e fadiga~\eqref{eq:updown}; com $b=5$, período de $1.1$~s (valor do modelo) e atividade em cerca de $35\%$ do tempo; com $b=1$, atividade contínua (fig.~\ref{fig:updown}) & Cálculo, script \texttt{models/sleep\_models.py}: período de $1.11$~s, fração de tempo ativo $0.35$ (média sobre um número inteiro de períodos); com $b=1$, fração ativa de $1.0$. O período do modelo é mais curto que o medido (linha~9) & --- & modelo \\
11 & O \nt{ACET} desliga a gangorra e leva o córtex ao funcionamento contínuo & A estimulação de um núcleo de neurônios \nt{ACET} em gato bloqueou o ritmo lento em $79\%$ dos neurônios do córtex e os levou a funcionamento contínuo, sobretudo pelo desaparecimento das hiperpolarizações longas. A acetilcolina suprime as correntes lentas de potássio que criam a fadiga & \cite{SteriadeAmzica1993,McCormick1992} & medido \\
12 & Os surtos de ritmo de $7$--$15$~Hz são gerados pela malha \nt{GLUT}--\nt{GABA} entre o córtex e um nó de retransmissão & Em fatias do nó de retransmissão visual de furão, os surtos são criados por rajadas de resposta das células de retransmissão à inibição vinda de um núcleo \nt{GABA} vizinho, que elas mesmas excitam & \cite{vonKrosigk1993,SteriadeMcCormickSejnowski1993} & medido \\
13 & Os replays do dia são acelerados: no hipocampo cerca de $20$ vezes, no córtex $6$--$7$ vezes & No sono de ondas lentas, as sequências repetidas de neurônios do hipocampo aparecem comprimidas no tempo. Depois de correr numa pista, as sequências de neurônios de lugar se repetiram na mesma ordem em surtos curtos de $\sim100$~ms, comprimidas cerca de $20$ vezes; no córtex, compressão de $6$--$7$ vezes & \cite{Nadasdy1999,LeeWilson2002,Euston2007} & medido \\
14 & Reforçar a coordenação dos replays com o córtex melhora a memória (relação causal) & Em ratos, após o treino, uma estimulação sincronizada com os replays reforçou sua ligação com as ondas lentas e os surtos de ritmo no córtex: no dia seguinte a memória era alta, e nos controles, no nível do acaso & \cite{Maingret2016} & medido \\
15 & A finalidade dos replays é o aprendizado intercalado da memória lenta & Teoria dos dois sistemas de aprendizado e cálculos em redes artificiais: o aprendizado sequencial estraga o antigo, o intercalado não. Não há experimento direto no cérebro mostrando que os replays servem justamente a isso & \cite{McClelland1995} & hipótese \\
16 & Após o sono, as sinapses diminuem proporcionalmente ao tamanho, e as grandes quase não mudam & Microscopia eletrônica tridimensional de $6920$ sinapses do córtex de camundongo: a área de contato após o sono é $\sim18\%$ menor que após a vigília, a redução é proporcional ao tamanho, e as sinapses grandes e estáveis não mudam & \cite{deVivo2017} & medido \\
17 & A tarefa principal do sono é a renormalização dos pesos & Hipótese da homeostase sináptica; as linhas 1 e 16 concordam com ela, mas não provam que seja a função principal & \cite{TononiCirelli2014} & hipótese \\
18 & O sono REM é um teste em bancada do modelo do mundo & O modo (tabela~\ref{tab:sleepmodes}) foi medido; que a rede execute o modelo do mundo é uma suposição teórica, sem experimento & \cite{Hobson2009} & hipótese \\
19 & Lavagem do cérebro no sono: questão em aberto & Um experimento: no sono, o espaço extracelular é $60\%$ mais largo e a limpeza, mais rápida. Outro, com outro método de medida: no sono e sob anestesia, a limpeza é bem mais lenta. Os resultados se contradizem & \cite{Xie2013,Miao2024} & hipótese \\
20 & Sono local: trechos do córtex de um animal acordado caem por instantes no silêncio & Em ratos, após vigília longa, os neurônios de um trecho isolado do córtex se calam brevemente, como no sono de ondas lentas, com uma onda lenta no EEG local; nesses momentos o rato erra mais na tarefa & \cite{Vyazovskiy2011} & medido \\
\bottomrule
\end{longtable}}
